% =============================================================================
% CONCLUSION AND RECOMMENDATIONS (Kết luận và Kiến nghị)
% =============================================================================
\chapter*{CONCLUSION AND RECOMMENDATIONS}
\phantomsection
\addcontentsline{toc}{chapter}{CONCLUSION AND RECOMMENDATIONS}

\section*{Conclusion}

Based on the analysis of 4,102 women aged 15--29 from the Vietnam MICS6 survey, this thesis demonstrates that the current paradigm of HPV vaccination in Vietnam is fundamentally inequitable.

\textbf{1. The Knowing--Doing Chasm:} While 62.4\% of young women are aware of the HPV vaccine, only 7.5\% have received it. This 55-percentage-point gap indicates that information diffusion, without concurrent financial accessibility, remains an insufficient public health strategy.

\textbf{2. Severe Pro-Rich Inequality:} The concentration index for vaccination (CI = 0.389) drastically eclipses that of awareness (CI = 0.149). Wealth does not merely influence vaccine uptake; it dictates it. 

\textbf{3. The Architecture of Disparity:} Erreygers decomposition isolates out-of-pocket cost as the ultimate barrier. Household wealth single-handedly drives 51.4\% of the inequality in vaccine uptake, compounded by mass media exclusion (23.8\%), rural isolation (17.1\%), and educational deficits (12.3\%).

\section*{Recommendations}

To eliminate cervical cancer, Vietnam must transition from opportunistic private vaccination to universal public provision:

\begin{enumerate}
	\item \textbf{Universal Public Financing.} The HPV vaccine must be immediately integrated into the national Expanded Programme on Immunization (EPI) with full state subsidies. Relying on private markets guarantees the systematic exclusion of the poorest quintiles.
	
	\item \textbf{Dismantling Geographic and Ethnic Barriers.} Conventional media campaigns fail remote populations. The Ministry of Health must deploy mobile vaccination units and village health workers to directly serve ethnic minority and rural communities, bypassing traditional health infrastructure bottlenecks.
	
	\item \textbf{School-Based Delivery Systems.} To bypass adult socioeconomic barriers entirely, vaccination must be universally administered through the public school system, guaranteeing equitable access prior to sexual debut.
	
	\item \textbf{Data-Driven Accountability.} Future longitudinal and provincial-level research is required to track whether subsequent policy interventions actually reduce the socioeconomic disparities mapped in this study, holding the health system accountable for equity.
\end{enumerate}

\clearpage
